\section{Consideraciones y Medidas de Seguridad para la Manipulación de Fibra Óptica}

\subsection{Advertencia General y Riesgos Asociados}
La manipulación de elementos relacionados con la fibra óptica presenta riesgos que dependen tanto de la naturaleza de la radiación lumínica como de las propiedades mecánicas del núcleo de vidrio. Para evitar accidentes personales y daños en los componentes del kit educativo, es obligatorio cumplir con las normas de seguridad durante el desarrollo de la práctica.

\subsection{Seguridad Óptica y Protección Ocular}
Las fuentes de luz utilizadas concentran alta densidad de potencia óptica en áreas reducidas. Dado que gran parte de las transmisiones operan en el espectro infrarrojo no visible, el ojo humano no activa el reflejo de parpadeo de forma natural.

\begin{itemize}
    \item \textbf{Prohibición de visión directa:} En ningún momento se debe ver directamente el extremo de una fibra óptica, conector o puerto emisor activo, ni utilizar instrumentos ópticos de aumento sin la debida verificación previa.
    
    \item \textbf{Direccionamiento del haz:} Queda estrictamente prohibido direccionar o apuntar el extremo de una fibra óptica o emisor hacia la cara o los ojos de otra persona.
    
    \item \textbf{Verificación de presencia de luz:} Para comprobar el paso de la señal óptica, se deben utilizar instrumentos auxiliares (como medidores de potencia óptica o fotodetectores) o reflejar el haz sobre superficies difusas a distancia segura, evitando el contacto ocular directo.
\end{itemize}

\subsection{Manipulación Mecánica y Peligro de Astillado}
Debido a que las fibras ópticas están constituidas por hilos de dióxido de silicio ($\text{SiO}_2$) cubiertos por capas protectoras, el corte o fractura inadecuada de estos hilos genera esquirlas de vidrio extremadamente transparentes, pequeñas y afiladas.

\begin{itemize}
    \item \textbf{Radio mínimo de curvatura:} Se deben evitar en todo momento curvas muy pronunciadas o estrangulamientos. Exceder el límite de curvatura puede quebrar o astillar la fibra en su interior, provocando pérdidas por radiación y el riesgo de cortes.
    
    \item \textbf{Prevención de heridas punzantes:} Las esquirlas de fibra pueden penetrar fácilmente en la piel o los ojos. No se deben tocar las puntas desnudas de la fibra ni recoger restos de vidrio con las manos sin protección.
    
    \item \textbf{Disposición de residuos:} Todos los fragmentos sobrantes de fibra cortada deben recogerse inmediatamente y depositarse en contenedores destinados a residuos de vidrio, evitando dejarlos sobre las mesadas de trabajo.
\end{itemize}

\subsection{Higiene, Limpieza y Manejo de Químicos}
Para mantener las terminales de los conectores libres de impurezas sin comprometer la seguridad del operador:

\begin{itemize}
    \item \textbf{Manipulación de Alcohol Isopropílico:} Debe emplearse alcohol isopropílico para la limpieza de conectores o fibras, lejos de fuentes de calor, chispas o arcos eléctricos, evitando el contacto prolongado con la piel y los ojos.
    
    \item \textbf{Higiene personal:} Al finalizar las actividades de laboratorio, o antes de frotarse los ojos, manipular lentes de contacto o ingerir alimentos, es indispensable lavarse bien las manos con abundante agua fría para retirar cualquier partícula adherida a la piel.
\end{itemize}